Storage is the only term in the groundwater flow equations that connects one time step to the next, so it is the term through which an adjoint solution propagates backward in time. This chapter describes that coupling and the selection between the specific-storage and specific-yield contributions.

\section*{Coupling between time steps}

The right-hand side at time step $k$ depends on the head at step $k-1$ through the storage terms. Differentiating the adjoint problem with respect to that earlier head gives the term that carries the later adjoint state backward, as in equation~\ref{eqn:carryback}. Writing it per cell,

\begin{equation}
	\label{eqn:drhsdh}
	\frac{\partial b}{\partial h^{k-1}} = -\left(
	\frac{\partial Q_{SS}}{\partial h^{k-1}} +
	\frac{\partial Q_{SY}}{\partial h^{k-1}} \right),
\end{equation}

\noindent where $Q_{SS}$ and $Q_{SY}$ are the volumetric flow rates from specific storage and specific yield.

The quantity in equation~\ref{eqn:drhsdh} is formed while the forward model runs and stored with that time step. It is paired with the adjoint state of the \emph{following} step, so the saturation it requires is the one at the step where it is formed, which the following step sees as the previous saturation.

\section*{Selecting the contributions}

\mf{} does not apply both contributions at once. The selection is made on the saturation of the cell at the previous time step, and it depends on the options in effect \citep{langevin2017}.

\subsection*{Specific yield}

Specific yield enters the previous-head derivative only through the saturation. For a cell that is neither full nor dry, $S_F^{k-1} = (h^{k-1} - BOT)/\Delta z$ and so $\partial S_F^{k-1} / \partial h^{k-1} = 1/\Delta z$, giving

\begin{equation}
	\label{eqn:dsy}
	\frac{\partial Q_{SY}}{\partial h^{k-1}} = \frac{S_y A}{\Delta t^{k}},
\end{equation}

\noindent where $S_y$ is the specific yield. Once the cell is full or dry the saturation stops changing with head, the derivative in equation~\ref{eqn:dsy} vanishes, and specific yield drops out.

\subsection*{Specific storage}

The specific-storage contribution depends on which formulation is active. Under the default formulation the previous-head derivative follows the previous saturation,

\begin{equation}
	\label{eqn:dss-default}
	\frac{\partial Q_{SS}}{\partial h^{k-1}} =
	\frac{SC1}{\Delta t^{k}} S_F^{k-1}, \qquad SC1 = S_s A \Delta z,
\end{equation}

\noindent where $S_s$ is the specific storage coefficient. Under the \texttt{SS\_CONFINED\_ONLY} option, specific storage acts only where the cell is full, and the term is removed rather than scaled:

\begin{equation}
	\label{eqn:dss-confined}
	\frac{\partial Q_{SS}}{\partial h^{k-1}} = \begin{cases}
		\dfrac{SC1}{\Delta t^{k}}, & S_F^{k-1} = 1, \\[2ex]
		0, & S_F^{k-1} < 1.
	\end{cases}
\end{equation}

A cell that is not convertible is treated as full at all times, and equation~\ref{eqn:dss-default} reduces to $SC1/\Delta t^{k}$.

Equations~\ref{eqn:dsy} through \ref{eqn:dss-confined} are complementary on the same test. Specific storage acts where the cell was full, specific yield where it was not.

\section*{Verification}

A seven by seven single-layer model was drawn down by a well over four time steps, leaving the cells between about half and four-fifths saturated. The sensitivity of the head at an observation cell to the pumping rate was compared with a finite-difference derivative taken by perturbing that rate.

\begin{table}[h]
	\centering
	\caption{Error in the head sensitivity against a finite-difference
	derivative, before and after the contributions were selected as described
	in this chapter.}
	\label{tab:sto-verify}
	\begin{tabular}{lrr}
		\toprule
		Formulation & Applying $SC1$ throughout & Selecting on saturation \\
		\midrule
		\texttt{SS\_CONFINED\_ONLY} & 10.7\% & 0.17\% \\
		Default                     & 2.6\%  & 0.15\% \\
		\bottomrule
	\end{tabular}
\end{table}

Applying $SC1$ at full cell thickness everywhere gives a partially saturated cell a storage term the forward model does not have, and Table~\ref{tab:sto-verify} shows the resulting error. The residual error after the correction is consistent with the nonlinearity of the finite-difference derivative itself.

The error is largest under \texttt{SS\_CONFINED\_ONLY} because that option removes the term entirely rather than reducing it, so the difference between applying and not applying it is the whole term. The size of the error in a given model scales with $S_s \Delta z$ relative to $S_y$: where a model sets $S_s$ such that $S_s \Delta z$ is comparable to or larger than $S_y$, an unselected specific-storage term dominates the coupling between time steps.

\section*{Limitation}

The \texttt{ORIGINAL\_SPECIFIC\_STORAGE} option, which restores the formulation used before \mf{} version 6.1.3, is not supported. Its previous-head derivative carries a term in the previous head itself in addition to the saturation, and \adj{} reports the option as unsupported rather than approximating it.
